\section{\texorpdfstring{$K$}{K}-代数的生成元}\label{sec:generators-of-k-algebras}

回顾一下, 我们一直要确定哪些环恰好可以作为坐标环. 上一节表明, 一个环 $R$ 要成为坐标环, 就必须是 $K$-代数. Proposition~\ref{prop:coordinate-ring-as-quotient-k-algebra} 进一步指出: $R$ 必须是某个多项式环的商. 但给定一个环, 它是否同构于这样的商未必一目了然. 例如, 在 Example~\ref{ex:complex-numbers-as-quotient-real-algebra} 中, 我们看到 $\mathbb{C}$ 作为 $\mathbb{R}$-代数同构于 $\mathbb{R}[x]$ 的一个商, 尽管它一开始并不是这样给出的.

本节的目标是准确刻画哪些 $K$-代数能作为多项式环的商. 为此所需的关键概念是多项式组合和 $K$-代数的生成元.

\begin{Definition}{多项式组合与生成元}\leavevmode\label{def:polynomial-combinations-and-generators}
  设 $A$ 是一个 $K$-代数, $S\subseteq A$ 是一个子集. $S$ 的一个\textbf{多项式组合}是 $A$ 中形如
  \[
    f(a_{1},\ldots,a_{n})
  \]
  的元素, 其中 $f\in K[x_{1},\ldots,x_{n}]$ 是某个多项式, 且 $a_{1},\ldots,a_{n}\in S$. $S$ 的全体多项式组合所成的集合称为 $A$ 中\textbf{由 $S$ 生成的子代数}, 记为 $K[S]$.
\end{Definition}

\BookPairedNote
  {读者不妨验证 $K[S]$ 确实是 $A$ 的子代数, 并且它是包含 $S$ 的最小子代数 (Exercise~\ref{ex:generated-subalgebra-is-smallest}).}
  {如果 $S=\{a_{1},\ldots,a_{n}\}$ 是有限集, 就省去集合的花括号, 写成 $K[S]=K[a_{1},\ldots,a_{n}]$.}
\noindent 另一种理解 $K[S]$ 的方式是: 它由 $A$ 中所有可以从 $S$ 中的元素出发, 通过加法, 乘法与用 $K$ 中的元素数乘而得到的元素组成.

\begin{Example}{{\textmd{$K[x,y]$} 的子代数}}\leavevmode\label{ex:subalgebras-of-two-variable-polynomials}
  考虑 $K$-代数 $A=K[x,y]$. 若令 $a=x$, 则由 $a$ 生成的子代数为 $x$ 的全体多项式:
  \[
    K[a]=K[x]\subseteq K[x,y].
  \]
  另一方面, 若 $b=x+y$, 则 $K[b]$ 是所有形如
  \[
    \sum_{i=0}^{k}r_{i}(x+y)^{i}
  \]
  的多项式所成的集合, 其中 $k\in\mathbb{N}$, 且对所有 $i$ 有 $r_{i}\in K$. 同时取生成元 $a$ 和 $b$, 可见 $K[a,b]$ 含有 $x=a$ 和 $y=b-a$. 从而可知 $K[a,b]$ 是 $x$ 和 $y$ 的全体多项式所成的集合:
  \[
    K[a,b]=K[x,y].
  \]
\end{Example}

\begin{Example}{{\textmd{$K[x,y]/\langle xy-1\rangle$} 的生成元}}\leavevmode\label{ex:generators-of-hyperbola-quotient-algebra}
  考虑 $K$-代数
  \[
    A=\frac{K[x,y]}{\langle xy-1\rangle}.
  \]
  $A$ 的任意元素都形如 $[f(x,y)]$, 其中 $f(x,y)\in K[x,y]$. 根据陪集运算的定义, 有
  \[
    [f(x,y)]=f([x],[y]).
  \]
  因此, $A$ 的任意元素都可以写成关于 $a=[x]$ 和 $b=[y]$ 的多项式表达式. 这意味着作为代数, $A$ 由 $a$ 和 $b$ 生成:
  \[
    A=K[a,b].
  \]
\end{Example}

虽然 Example~\ref{ex:generators-of-hyperbola-quotient-algebra} 中的记号 $K[a,b]$ 让人想起多项式环, 但 $A$ 并不等于以 $a$ 和 $b$ 为变量的多项式环. 特别地, 生成元之间有一个关系:
\[
  ab=1.
\]
在这里以及全书中, 我们用字母表末尾的字母 (例如 $x$ 和 $y$) 表示多项式环中的变量. 根据定义, 这些变量之间没有任何关系, 也就是说, 两个多项式相等当且仅当它们的系数相同. 另一方面, 我们用字母表开头的字母 (例如 $a$ 和 $b$) 表示 $K$-代数的生成元, 它们可以满足一些关系. 换句话说, 即使 $f$ 与 $g$ 是不同的多项式, 也可能有 $f(a,b)=g(a,b)$.

前面两个例子中的整个 $K$-代数都可由有限多个元素生成. 为了描述这种性质, 我们称它们为有限生成的 $K$-代数.

\begin{Definition}{有限生成的 $K$-代数}\leavevmode\label{def:finitely-generated-k-algebra}
  设 $A$ 是一个 $K$-代数. 如果存在 $a_{1},\ldots,a_{n}\in A$, 使得 $A=K[a_{1},\ldots,a_{n}]$, 就称 $A$ 是\textbf{有限生成的}.
\end{Definition}

例如, 多项式环 $K[x_{1},\ldots,x_{n}]$ 是有限生成的 $K$-代数, 只须对所有 $i=1,\ldots,n$ 取 $a_{i}=x_{i}$. 更一般地, 对任意理想 $I$, 商环 $K[x_{1},\ldots,x_{n}]/I$ 都是有限生成的, 因为可对所有 $i=1,\ldots,n$ 取 $a_{i}=[x_{i}]$, 这是 Example~\ref{ex:generators-of-hyperbola-quotient-algebra} 的推广. 实际上, 在同构意义下, 这些是有限生成的 $K$-代数的全部例子. 下面验证这一点.

\begin{Proposition}{有限生成的 $K$-代数的刻画}\leavevmode\label{prop:characterization-of-finitely-generated-k-algebras}
  设 $A$ 是一个 $K$-代数. 则 $A$ 是有限生成的, 当且仅当对某个 $n\in\mathbb{N}$ 和某个理想 $I\subseteq K[x_{1},\ldots,x_{n}]$, 有 $K$-代数同构
  \[
    A\cong K[x_{1},\ldots,x_{n}]/I.
  \]
\end{Proposition}

\begin{Proof}
  假设 $A$ 是有限生成的. 根据定义, 这意味着存在 $a_{1},\ldots,a_{n}\in A$, 使得 $A=K[a_{1},\ldots,a_{n}]$. 定义函数
  \[
    \begin{aligned}
      \phi:K[x_{1},\ldots,x_{n}]&\to A\\
      f(x_{1},\ldots,x_{n})&\mapsto f(a_{1},\ldots,a_{n}).
    \end{aligned}
  \]
  容易验证 $\phi$ 是 $K$-代数同态, 而 $A=K[a_{1},\ldots,a_{n}]$ 等价于 $\phi$ 是满射. 令 $I=\ker(\phi)$, 由第一同构定理可知
  \[
    A\cong K[x_{1},\ldots,x_{n}]/I.
  \]

  反过来, 假设存在同构
  \[
    \psi:K[x_{1},\ldots,x_{n}]/I\to A,
  \]
  并对 $i=1,\ldots,n$ 令 $a_{i}=\psi([x_{i}])$. 我们要证明 $A=K[a_{1},\ldots,a_{n}]$. 假设 $a\in A$, 必须证明 $a$ 是关于 $a_{1},\ldots,a_{n}$ 的一个多项式表达式. 因为 $\psi$ 是满射, 存在 $[f]\in K[x_{1},\ldots,x_{n}]/I$, 使 $\psi([f])=a$. 那么
  \[
    \begin{aligned}
      a&=\psi([f(x_{1},\ldots,x_{n})])=\psi(f([x_{1}],\ldots,[x_{n}]))\\
       &=f(\psi([x_{1}]),\ldots,\psi([x_{n}])),
    \end{aligned}
  \]
  其中第二个等号来自陪集运算, 第三个等号来自 $\psi$ 为 $K$-代数同态的假设. 因为 $a_{i}=\psi([x_{i}])$, 可见 $a$ 是关于 $a_{1},\ldots,a_{n}$ 的多项式表达式, 因而 $A$ 是有限生成的.
\end{Proof}

Proposition~\ref{prop:characterization-of-finitely-generated-k-algebras} 中这种形式的同构常称为 $K$-代数 $A$ 的一个\textbf{表示}. $x_{1},\ldots,x_{n}$ 的像是这个表示的生成元, $I$ 中的多项式是这个表示的关系.

\BookPairedNote
  {由 Proposition~\ref{prop:coordinate-ring-as-quotient-k-algebra}, 仿射簇的坐标环典范同构于多项式环的商. 因此, 由 Proposition~\ref{prop:characterization-of-finitely-generated-k-algebras} 得到下面的结果.}
  {关于几个不是有限生成的代数的例子, 见 Exercises~\ref{ex:infinite-variable-polynomial-algebra-not-finite}--\ref{ex:subalgebra-not-finitely-generated}.}

\begin{Corollary}{{\textmd{$K[X]$} 是有限生成的}}\leavevmode\label{cor:coordinate-ring-finitely-generated}
  设 $X\subseteq\mathbb{A}^{n}$ 是仿射簇. 则坐标环 $K[X]$ 是有限生成的 $K$-代数.
\end{Corollary}

更明确地说, Proposition~\ref{prop:characterization-of-finitely-generated-k-algebras} 的证明表明, $K[X]$ 作为 $K$-代数, 可以由 Proposition~\ref{prop:coordinate-ring-as-quotient-k-algebra} 的典范同构
\[
  K[x_{1},\ldots,x_{n}]/\mathcal{I}(X)=K[X]
\]
中的元素 $[x_{1}],\ldots,[x_{n}]$ 生成. 这些像就是 $n$ 个坐标函数
\[
  \begin{aligned}
    C_{i}:X&\to K\\
    (a_{1},\ldots,a_{n})&\mapsto a_{i}.
  \end{aligned}
\]
因此, $K[X]$ 是由 $X$ 上的坐标函数生成的 $K$-代数. 这恰好就是我们把它称为 $X$ 的\textquotedbl{}坐标环\textquotedbl{}的原因.

\BookExerciseHeading
\begin{BookExercises}
  \item\label{ex:generated-subalgebra-is-smallest} 设 $A$ 是一个 $K$-代数, $S\subseteq A$ 是一个子集.
    \begin{enumerate}[label=(\alph*)]
      \item 证明 $K[S]$ 是 $A$ 的子代数.
      \item 证明 $K[S]$ 包含在 $A$ 的每个包含 $S$ 的子代数中.
    \end{enumerate}
  \item\label{ex:two-generating-sets-of-polynomial-algebra} 给出三元素生成 $K[x,y,z]$ 的两个不同例子.
  \item\label{ex:one-generator-for-parabola-quotient} 考虑 $K$-代数
    \[
      A=K[x,y,z]/\langle y-x^{2},z\rangle.
    \]
    找出一个元素 $a\in A$, 使 $A=K[a]$.
  \item\label{ex:hyperbola-quotient-needs-two-generators} 考虑 $K$-代数
    \[
      A=K[x,y]/\langle xy-1\rangle.
    \]
    证明对任意 $a\in A$, 都有 $A\ne K[a]$.
  \item\label{ex:algebra-maps-agree-on-generators} 假设 $A$ 与 $B$ 是 $K$-代数, 且对某些 $a_{1},\ldots,a_{n}\in A$, 有 $A=K[a_{1},\ldots,a_{n}]$. 证明如果 $\phi,\psi:A\to B$ 是 $K$-代数同态, 且对所有 $i=1,\ldots,n$ 都有 $\phi(a_{i})=\psi(a_{i})$, 则对所有 $a\in A$ 都有 $\phi(a)=\psi(a)$.
  \item\label{ex:presentations-of-rational-subalgebras} 把 $\mathbb{R}$ 视为 $\mathbb{Q}$-代数.
    \begin{enumerate}[label=(\alph*),ref=(\alph*)]
      \item\label{item:present-quadratic-rational-algebra} 子代数 $\mathbb{Q}[\sqrt{2}]\subseteq\mathbb{R}$ 是有限生成的, 所以由 Proposition~\ref{prop:characterization-of-finitely-generated-k-algebras} 可知, 对某个 $n$ 和某个 $I$, 有 $\mathbb{Q}[\sqrt{2}]\cong\mathbb{Q}[x_{1},\ldots,x_{n}]/I$. 明确找出满足此式的一个 $n$ 和一个 $I$.
      \item 对子代数 $\mathbb{Q}[\sqrt[3]{2}]\subseteq\mathbb{R}$ 重做 \ref{item:present-quadratic-rational-algebra}.
      \item 对子代数 $\mathbb{Q}[\pi]\subseteq\mathbb{R}$ 重做 \ref{item:present-quadratic-rational-algebra}.
    \end{enumerate}
  \item\label{ex:infinite-variable-polynomial-algebra-not-finite} 证明 $K[x_{1},x_{2},x_{3},\ldots]$ 不是有限生成的 $K$-代数. (环 $K[x_{1},x_{2},x_{3},\ldots]$ 在 Example~\ref{ex:infinitely-generated-variable-ideal} 中定义.)
  \item\label{ex:finite-type-over-q-countable} 证明任何有限生成的 $\mathbb{Q}$-代数都是可数的, 并据此得出 $\mathbb{R}$ 不是有限生成的 $\mathbb{Q}$-代数.
  \item\label{ex:subalgebra-not-finitely-generated} 考虑子代数
    \[
      A=K[x,xy,xy^{2},xy^{3},xy^{4},\ldots]\subseteq K[x,y].
    \]
    证明 $A$ 不是有限生成的 $K$-代数. 因此, 有限生成代数的子代数不一定是有限生成的.
\end{BookExercises}
